% Hand-written closing section.  Nothing generates this file.  It is the last
% body section: \input before the acknowledgments in main.tex, so its \section
% closes the rule opened by the final case.  The counts in the first sentence
% are the only numbers here, and they come from the case.json files
% (14 cases, 18 results, 16 exact + 2 computer-assisted); `make check` will not
% catch it if they drift, so re-count before changing them.
\section{Conclusion}

\noindent The archive behind this manuscript currently holds eighteen refuted statements across fourteen cases; sixteen of them reduce to exact rational arithmetic and two to rigorous interval enclosures, and every one recomputes from source in the accompanying repository (several of the refutations are in additional supported by an analytic proof too, not just numerical). %That is the only claim we make for the collection as a whole. Each statement is quoted as it was posed, with its hypotheses and quantifiers left alone, and each witness is small enough to be checked --- by hand where we could manage it, by machine everywhere.

% \vskip5pt
% \noindent This is also why we think the asymmetry is worth pressing. A model's transcript is not evidence and we do not ask the reader to treat it as such; but a counterexample is a finite object, and a finite object either survives the checks or it does not. The search may be as heuristic as one likes, since nothing of it is load-bearing by the time the certificate is written.

\vskip5pt
\noindent What we find most interesting, though, is that the better counterexamples say more than ``false.'' They locate a boundary. The quantum coupon-collector sum is positive definite through five variables in every dimension and fails at six; the mixed-norm inequality holds at every positive integer exponent and fails at every non-integer one; and in the same rank-two instance where Schur convexity of the $t^\delta$-normalized Macdonald ratio fails, the $1^n$-normalized ratio satisfies the very inequality violated. A refutation of that kind does not close a subject so much as move it, and the corrected question it leaves behind is usually the one worth asking next.

%%% Local Variables:
%%% mode: LaTeX
%%% TeX-master: "main"
%%% End:
